% =============================================================================
% Section 2 -- Background and Related Work
% =============================================================================
\section{Background and Related Work}
\label{sec:background}

\subsection{Agent Evaluation and the Deployment Gap}
\label{subsec:deployment-gap}

Pre-deployment evaluation of AI systems has matured substantially. The
emergence of standardized benchmarks~\cite{liang2022holistic}, automated
red-teaming pipelines~\cite{ganguli2022redteaming,perez2022redteaming},
behavioral testing frameworks such as CheckList~\cite{ribeiro2020checklist},
and alignment techniques such as RLHF~\cite{ouyang2022instructions} and
Constitutional AI~\cite{anthropic2023constitutional} have created a well-tooled
pre-release lifecycle. Post-deployment monitoring has received comparatively
less systematic attention. Pan, Arabzadeh, and Ellis~\cite{pan2025measuring}
conducted the first large-scale empirical study of production AI agents,
finding that most deployments fail or underdeliver and that little is publicly
known about post-release management. Sapra et al.~\cite{sapra2025agentcompass}
address detection via AgentCompass---structured post-deployment
debugging---but prescribe no remediation. Rath~\cite{rath2026drift} formalizes
agent drift across 12 measurable dimensions, projecting a 42\% task success
reduction for unchecked drift via simulation. The APIP addresses what all three
leave open: a structured remediation protocol once degradation is identified.

Debugging toolkits have recently begun to close a detection-to-recovery loop at
the level of a single failed execution. AgentDebugX~\cite{zhu2026agentdebugx}
organizes debugging as a cycle of detection, attribution, recovery, and rerun,
verifies each repair by re-executing the failed trajectory, and maintains a
shared corpus of diagnosed failures. The APIP operates one level above such
tooling, and the two compose. A successful rerun verifies an instance; in the
terms of Section~\ref{subsec:kirkpatrick}, this is Level 1/2 evidence rather
than evidence of durable recovery against the behavioral contract. A run-level
loop also does not decide whether an aggregate degradation warrants intervention
at all, when to escalate intervention tiers, or when an agent should be retired
rather than repaired, and it produces no governance record of those decisions.
Within an APIP, repair tooling of this kind is a natural executor for Tier 1 and
Tier 2 interventions during the remediation window, while the APIP supplies the
trigger discipline of the behavioral contract, the bounded window with formal
exit evaluation, and the audit trail required by the frameworks in
Section~\ref{subsec:governance}.

\subsection{The Performance Improvement Plan as a Structural Template}
\label{subsec:pip-template}

The Performance Improvement Plan in organizational management is a formal
document issued to an underperforming employee that specifies: the performance
deficiencies observed, the root causes identified, the corrective actions
required, the time window for improvement, the check-in schedule, and the
criteria by which success or failure will be evaluated. The American Society
for Human Resource Management recommends that a PIP should be used when there
is a genuine commitment to help the employee improve, not merely as a precursor
to termination~\cite{kirkpatrick2016four}.
Kirkpatrick~\cite{kirkpatrick1994evaluating,kirkpatrick2016four} emphasizes
that the PIP should be developed collaboratively by the manager and employee,
because remediation requires both parties' participation to succeed.

The PIP analogy to agent management is structurally productive but requires
important qualification. A human employee can engage in self-directed
improvement: they understand the feedback, adjust their behavior, and exercise
agency. A deployed AI agent cannot. The APIP is therefore a management protocol
imposed entirely by human operators---the agent is the subject, not the author,
of the plan. This distinction is theoretically important and has practical
implications for where we locate the locus of accountability in the APIP
lifecycle. In organizational terms, the APIP collapses the manager--employee
dyad: the operator is both the manager who writes the plan and the agent of
change who implements it.

\subsection{Kirkpatrick's Four-Level Evaluation Framework}
\label{subsec:kirkpatrick}

Kirkpatrick's~\cite{kirkpatrick1994evaluating} four-level training evaluation
model, first published in 1959 and refined over the subsequent six decades,
provides a hierarchical framework for evaluating the impact of interventions:
Level 1 (Reaction) assesses immediate participant response, Level 2 (Learning)
measures knowledge or skill acquisition, Level 3 (Behavior) evaluates whether
learning transfers to on-the-job behavior change, and Level 4 (Results)
measures organizational outcomes~\cite{kirkpatrick1994evaluating,kirkpatrick2016four}.
The model has been applied across business, government, military, and
healthcare contexts, and remains the most widely used training evaluation
framework globally.

We adopt the Kirkpatrick framework as a structural template for evaluating APIP
interventions. When a prompt-level, retrieval-level, or model-level
intervention is applied to a degraded agent, the check-in evaluation should
assess impact at multiple levels: Did the intervention change the agent's
immediate output patterns (Level 1/2)? Did it produce durable behavioral change
in how the agent handles the failure-mode class (Level 3)? Did it improve the
business metrics specified in the behavioral contract---CSAT, task completion,
escalation rate (Level 4)? An intervention that improves Level 1/2 metrics
without producing Level 3/4 improvement has not remediated the agent; it has
merely suppressed the symptom. This hierarchical evaluation discipline is
essential for distinguishing genuine recovery from surface-level metric
manipulation.

\subsection{AI Governance and Documentation Requirements}
\label{subsec:governance}

Regulatory frameworks for AI systems are increasingly emphasizing
documentation, human oversight, and traceability. The EU AI
Act~\cite{eu2024aiact} requires providers of high-risk AI systems to implement
post-market monitoring systems and maintain logs sufficient to identify sources
of error. US executive orders on AI safety have similarly emphasized the
importance of human oversight in consequential automated decision systems.
Brundage et al.~\cite{brundage2020trustworthy} argue that verifiable claims
require mechanisms beyond pre-deployment testing alone. At FAccT 2024, Kolt,
Heim, and Anderljung~\cite{kolt2024visibility} identified three visibility
mechanisms for AI agents---identifiers, real-time monitoring, and activity
logging---arguing accountability requires linking actions to specific agents and
operators. The APIP schema's \texttt{apip\_id}, \texttt{agent\_id},
\texttt{owner}, and \texttt{exit\_evaluated\_at} fields operationalize this
linkage, positioning the APIP as both a technical practice and a governance
instrument.

\subsection{AI Control Protocols and the Threat-Model Boundary}
\label{subsec:ai-control}

A complementary line of work, the AI Control literature originating with
Greenblatt et al.~\cite{greenblatt2023control}, addresses post-deployment
safety under an explicitly adversarial threat model. Control protocols assume
the deployed model may be scheming and design runtime safeguards---trusted
monitoring (a weaker, trusted model audits an untrusted model's actions),
untrusted monitoring (one instance of the untrusted model monitors another),
defer-to-trusted, and resampling protocols---intended to maintain safety even
when the agent is adversarial. Subsequent work has surfaced failure modes of
these protocols, including capability gaps in trusted
monitors~\cite{bhatt2025ctrlz}, collusion between untrusted monitor and
policy~\cite{mallen2025subversion}, and adaptive prompt-injection attacks that
evade diverse monitors and can be amplified by
resampling~\cite{schoen2025stress}.

The APIP framework is complementary to AI Control rather than competing with
it. The two operate at different threat models, granularities, and timescales.
Control protocols address the question ``is this individual action safe given a
possibly-scheming agent?'' through per-action runtime decisions. The APIP
addresses the question ``is this agent meeting its behavioral contract over an
aggregated time window?'' through structured remediation across days to weeks.
Control protocols presuppose adversarial intent; APIP's primary failure modes
(Section~\ref{sec:taxonomy}) presuppose capability degradation or value drift.
The two regimes also consume and produce different signals: control protocols
generate per-action audit decisions that can serve as input to APIP triggers (a
sustained rise in audit-flagged actions is itself a behavioral-contract
breach), while APIP's aggregated cause attribution can identify when a
degradation pattern warrants escalation to control-protocol-level intervention.
We position the \texttt{adversarial\_subversion} failure mode
(Section~\ref{sec:taxonomy}) as the integration point between the two regimes:
when cause attribution identifies an adversarial pattern rather than capability
drift, the appropriate intervention is drawn from the control literature rather
than from APIP's standard tier model.

\subsection{Team Dynamics, Membership Change, and Coordination}
\label{subsec:team-dynamics}

A growing body of organizational behavior research addresses how team
composition changes affect team functioning and performance. Morgeson,
Mitchell, and Liu~\cite{morgeson2015event} introduced Event System Theory
(EST), which posits that organizational events are characterized by three
properties---novelty, disruption, and criticality---that determine their impact
on organizational entities. Events that are higher on these dimensions are more
likely to produce changes in behavior, features, or subsequent events.

Summers, Humphrey, and Ferris~\cite{summers2012team} applied an event-based
lens to team member change and found that membership changes caused high levels
of flux in coordination when the change involved a strategically core role or
when information transfer during the transition was low. Trainer et
al.~\cite{trainer2020membership} reconceptualized each team membership change as
a discrete team-level event with varying degrees of novelty, disruptiveness,
and criticality, finding that newcomer entry can substantially impact team
functioning beyond its impact on the newcomer themselves.

Research on Transactive Memory Systems (TMS)---collective memory systems in
which team members specialize in different knowledge domains while maintaining
shared awareness of who knows
what~\cite{wegner1987transactive,lewis2007group,ren2011transactive,liang1995group,moreland2003transactive}---provides
a cognitive mechanism for understanding how membership changes disrupt
coordination. When a new member enters, the team's TMS becomes inaccurate:
existing members' calibrations of who knows what and how to coordinate are
invalidated~\cite{lewis2007group,wegner1987transactive}. The newcomer
socialization literature further documents that support for newcomers declines
within the first 90 days~\cite{kammeyer2013support}, and that structured
onboarding significantly improves adjustment outcomes~\cite{bauer2025newcomer}.

These streams provide the theoretical vocabulary for
Section~\ref{subsec:mesh}: how a new agent introduced to an existing mesh can
degrade peer performance even when performing its own task well.
